\section{Method}
\label{sec:method}

\begin{definition}[Validated edit]
A validated edit is a tuple
$e=(f,b_0,b_1,r,h)$ containing a file identifier, half-open byte interval,
replacement text, and the expected hash of the replaced bytes.
\end{definition}

Let $S_f$ be the authoritative bytes for file $f$. An edit is eligible exactly
when its interval is valid and its precondition matches:
\begin{equation}
  \label{eq:eligibility}
  0 \le b_0 \le b_1 \le \lvert S_f \rvert
  \quad\land\quad
  H(S_f[b_0:b_1]) = h.
\end{equation}

\begin{theorem}[Stale-edit exclusion]
If the bytes inside an edit's interval change after the edit is proposed, a
collision-resistant expected-slice hash causes validation to reject the edit,
except with negligible collision probability.
\end{theorem}

\begin{proof}
Changing the interval changes the hash input. Acceptance would therefore
require a collision or second preimage for the selected digest.
\end{proof}

The benchmark applies an edit, reparses the candidate, checks complete source
coverage, and then applies the inverse edit. A compact representation is:

\begin{lstlisting}[language={},caption={Illustrative edit loop.},label={lst:loop}]
for edit in edits:
    candidate = validate_and_apply(source, edit)
    assert complete_coverage(candidate)
    source = apply_inverse(candidate, edit)
\end{lstlisting}
